%%%%%%%%%%%%%%%%%%%%%%%%%%%%%%%%%%%%%%%%%%%%%%%%%%%%%%%%%%%%%%%%%%%%%%%%%%%%%%%
% Frontiers in Artificial Intelligence -- Revised Original Research Article
% This v2 preserves paper/v1 unchanged and reflects the completed experiments.
%%%%%%%%%%%%%%%%%%%%%%%%%%%%%%%%%%%%%%%%%%%%%%%%%%%%%%%%%%%%%%%%%%%%%%%%%%%%%%%

\documentclass[utf8]{FrontiersinHarvard}

\usepackage{url,hyperref,lineno,microtype}
\usepackage[singlespacing]{setspace}
\usepackage{amsmath}
\usepackage{booktabs}
\usepackage{graphicx}

\linenumbers

\def\keyFont{\fontsize{8}{11}\helveticabold}
\def\firstAuthorLast{Rhrissi {et~al.}}
\def\Authors{Zhour Rhrissi\,$^{1,*}$, Sanaa El Fkihi\,$^{1}$ and Rdouan Faizi\,$^{1}$}
\def\Address{$^{1}$National Higher School for Computer Science and Systems Analysis (ENSIAS), Mohammed V University in Rabat, Rabat, Morocco}
\def\corrAuthor{Zhour Rhrissi}
\def\corrEmail{zhour\_rhrissi@um5.ac.ma}

\begin{document}
\onecolumn
\firstpage{1}

\title[Cross-generator review detection]{What Transfers Across Generators? Efficient Detection of Machine-Generated Product Reviews with a 0.5B-Parameter Model}

\author[\firstAuthorLast]{\Authors}
\address{\Address}
\correspondance{Correspondence: \corrAuthor, \corrEmail}
\extraAuth{}

\maketitle

\begin{abstract}
\section{}
High accuracy on a single machine-generated-text dataset does not establish that a detector recognizes machine authorship rather than one generator's fingerprint. We study this distinction for English product reviews generated by four model families: a historical GPT-2 corpus and three modern generators recorded as LFM2.5-1.2B, DeepSeek-v4-pro, and GLM-5.2. A Qwen2.5-0.5B sequence classifier was adapted with rank-stabilized quantized low-rank adaptation (Rs-QLoRA). Generator-specific experiments used three training seeds; a balanced leave-one-generator-out (LOGO) evaluation and a human-only permutation control tested transfer and leakage. Matched-distribution accuracy was 98.50--99.78\%. Transfer, however, was structured: mean off-diagonal recall was 39.3\% across all four generators, but 78.4\% among the three modern generators. Detectors trained on DeepSeek or GLM outputs transferred broadly to the other modern generators, whereas the LFM-trained boundary was narrower; transfer between GPT-2 and every modern generator was at most 0.5\%. In LOGO evaluation, recall on an unseen modern generator was 91.1--96.7\% at a 1.92--2.13\% human false-positive rate, while GPT-2 recall was 0.1\%. This LOGO result holds out generator identity but not prompt-source identity: 94.3--94.5\% of modern target source reviews also occur through another generator in pooled training. A random-label human-only control produced chance performance. Against standard QLoRA on one matched modern task, Rs-QLoRA had similar accuracy but reduced mean training time from 32.0 to 17.3 minutes on one H100 GPU. The evidence supports an efficient detector of transferable signals within the tested modern generation pipeline, not a universal detector of machine-written reviews.

\tiny
\keyFont{\section{Keywords:} machine-generated text detection, product reviews, cross-generator transfer, Qwen2.5, parameter-efficient fine-tuning, Rs-QLoRA, generator shift}
\end{abstract}

\section{Introduction}

Online reviews affect consumer decisions and platform markets, making review manipulation a longstanding concern \citep{Jindal2008,Ott2011}. The emergence of fluent text generators adds a related but distinct problem: determining whether a review was written by a person or generated by a model. Machine authorship is not equivalent to deception. A generated review can be disclosed and benign, while a human review can be deceptive. We therefore use \emph{human-authored} (OR) and \emph{computer-generated} (CG) as the class labels, reserving \emph{fake} or \emph{deceptive} for contexts in which intent has been established.

A common evaluation trains and tests a detector on outputs from the same generator. Such a test measures in-distribution discrimination, but its score can be dominated by generator-specific lexical, decoding, or data-pipeline artifacts. Prior work on machine-generated text detection shows that detectors can fail under unseen generators, domains, decoding strategies, and adversarial transformations \citep{Pu2023,Dugan2024}. Consequently, a near-ceiling matched score is not evidence of a universal machine-text signal.

This study asks three narrower questions. \textbf{RQ1:} How much does matched-distribution performance overstate cross-generator transfer? \textbf{RQ2:} Is transfer structured across historical and modern generators, and is it directional? \textbf{RQ3:} What does rank stabilization contribute when the base model is quantized and adapted for this task? We evaluate a 0.5B-parameter Qwen2.5 classifier across a four-generator matrix, a generator-held-out pooled setting, matched classical and transformer baselines, an Rs-QLoRA versus standard QLoRA comparison, and a human-only random-label control.

The principal contribution is not another isolated accuracy record. It is an empirical separation between matched detection and transfer. We find a sharp GPT-2--modern boundary and strong but asymmetric transfer among the tested modern generators. We also identify a key boundary on this result: the modern corpora share a grounded prompting protocol and source-review pool. Thus, the experiments establish transfer under controlled source and prompting support, while a joint generator-and-source holdout remains necessary before claiming broader out-of-distribution generalization.

\section{Related Work}

\subsection{Deceptive reviews and machine-generated reviews}

Early opinion-spam research studied duplicate, anomalous, or intentionally deceptive reviews \citep{Jindal2008}. The Ott corpus supplied a gold-standard task involving deliberately deceptive hotel reviews and demonstrated that linguistic classifiers could separate them from truthful reviews \citep{Ott2011}. These tasks concern veracity or intent, neither of which follows from authorship alone.

The dataset underlying the historical portion of our study was introduced by \citet{Salminen2022}, who used GPT-2 to create product reviews paired with human reviews and evaluated automatic detection. We retain that corpus as a historical generator domain, but we do not treat its CG label as proof of deception. Our modern corpora update the generator side while preserving the product-review domain and the OR source population.

\subsection{Generalization of machine-generated text detectors}

Cross-generator robustness has become a central concern in machine-generated text detection. \citet{Pu2023} trained detectors on outputs from individual generators and found structured but incomplete zero-shot transfer, including benefits from mixing generator sources. RAID expanded evaluation across models, domains, decoding strategies, and attacks, showing that high headline accuracy can collapse under realistic distribution shifts \citep{Dugan2024}. These findings motivate our use of a full train-generator by test-generator matrix and LOGO training rather than relying only on matched tests.

Our setting is intentionally narrower than general-purpose benchmarks: English product reviews with aligned categories, ratings, and source-review grounding. This control reduces domain shift but introduces shared-prompt and source dependence. We measure and report that dependence explicitly.

\subsection{Parameter-efficient adaptation}

LoRA freezes a pretrained model and learns low-rank updates to selected weight matrices \citep{Hu2022}. QLoRA further backpropagates through a quantized frozen base into LoRA adapters, substantially reducing adaptation memory \citep{Dettmers2023}. Standard LoRA scales the adapter contribution by $\alpha/r$. Rank-stabilized LoRA instead uses $\alpha/\sqrt{r}$ to avoid increasingly weak updates at larger ranks \citep{Kalajdzievski2023}. We combine a 4-bit Qwen2.5 base \citep{Qwen2024} with rank-64 adapters and compare rank-stabilized and standard scaling under identical modern-data and hardware conditions.

\section{Data and Experimental Design}

\subsection{Human and generated reviews}

All datasets are balanced binary corpora with ten product categories and fields for category, rating, label, and text. OR reviews originate from the public product-review corpus used by \citet{Salminen2022}. Its historical GPT-2 CG half forms the first generator domain. For each modern domain, the human half comes from the same OR pool and the CG half is newly generated. Table~\ref{tab:data} summarizes the final deduplicated datasets.

\begin{table}[h!]
\centering
\caption{Final datasets after normalization, exact deduplication, and class balancing. Model names for modern generators are the request identifiers recorded at generation time; provider-side immutable revisions were not retained.}
\label{tab:data}
\begin{tabular}{lrrr}
\toprule
Generator domain & OR & CG & Total \\
\midrule
GPT-2 historical & 20,190 & 20,190 & 40,380 \\
LFM2.5-1.2B & 20,212 & 20,212 & 40,424 \\
DeepSeek-v4-pro & 20,212 & 20,212 & 40,424 \\
GLM-5.2 & 20,208 & 20,208 & 40,416 \\
\bottomrule
\end{tabular}
\end{table}

Each modern CG review was grounded on one OR review. The generator received its category, integer rating, the first 400 characters of the source review, and a target length of 4--120 words. It was instructed to write a new review for the same product while emphasizing one of seven deterministically selected angles. Generation used up to five attempts, temperatures 0.85--1.09, $top\_p=0.95$, and a word-set Jaccard acceptance threshold below 0.6 relative to the source. No API decoding seed was supplied. Unicode punctuation was normalized for both classes, empty outputs and exact within-CG duplicates were removed, and any CG text exactly matching an OR text after case folding was excluded. The exact prompts, acceptance rules, artifact hashes, provider endpoints, and known unknowns are recorded in the public provenance manifest.

The modern generators share this prompt template, acceptance logic, product domain, and OR source pool. These controls make generator comparisons less confounded by topic and rating, but they can also create a common pipeline signature. Generator identity is therefore inseparable from provider version, prompting, decoding, and corpus construction in the current study.

\subsection{Frozen partitions and source dependence}

Each generator dataset used a fixed 80\%/6\%/14\% train/validation/test partition produced with split seed 1998. Training seeds never changed the partition. These were row-level, not source-grouped, splits. Exact-text deduplication prevents identical reviews from crossing partitions, but a generated review and the OR review that grounded it may occur in different partitions.

The cross-generator design adds a second dependence. Modern generators were prompted from the same source-review pool, so the same source ID may underlie a training generation from one model and a test generation from another. We recovered source IDs from the raw generation records and audited this overlap (Table~\ref{tab:overlap}). Pairwise overlap covers 69.3--81.5\% of target sources. In pooled LOGO data it covers 94.3--94.5\%. Accordingly, LOGO holds out generator identity and excludes held-out OR texts from pooled training, but it does not hold out source-prompt identity. We use \emph{generator-held-out transfer} rather than the stronger \emph{fully out-of-distribution generalization} for this evaluation.

\begin{table}[h!]
\centering
\caption{Fraction of a modern target's CG test source IDs also represented by CG training examples from another generator. Opposite transfer directions have similar overlap, which is useful when interpreting directional recall differences.}
\label{tab:overlap}
\begin{tabular}{lccc}
\toprule
Training source & LFM target & DeepSeek target & GLM target \\
\midrule
LFM & -- & 69.3\% & 80.8\% \\
DeepSeek & 69.8\% & -- & 80.2\% \\
GLM & 81.5\% & 80.8\% & -- \\
Pooled other modern generators & 94.5\% & 94.5\% & 94.3\% \\
\bottomrule
\end{tabular}
\end{table}

\subsection{Evaluation protocols}

We use five complementary protocols.

\begin{enumerate}
\item \textbf{Matched:} train and test within each generator domain using seeds 1998, 7, and 42.
\item \textbf{Cross-generator matrix:} evaluate each generator-specific detector on the CG test examples from every generator. We report CG recall because the target question is whether unseen generator outputs are caught. Each detector's human false-positive rate (FPR) is reported separately to prevent a flag-everything interpretation.
\item \textbf{LOGO:} pool budget-matched training data from the other three generators and test on the held-out generator. Duplicate OR texts are collapsed, and every exact OR text in the target test partition is excluded from pooled data. CG examples are sampled equally from the training generators. LOGO neural results currently use one training seed (1998).
\item \textbf{Human-only permutation control:} randomly assign balanced pseudo-labels to 20,214 OR reviews and train the same detector. Above-chance performance would indicate label or pipeline leakage unrelated to machine authorship.
\item \textbf{Matched baselines and adaptation comparison:} compare TF--IDF, three fully fine-tuned encoders, and standard versus rank-stabilized QLoRA on the frozen DeepSeek split.
\end{enumerate}

The primary metrics are accuracy, CG recall, human FPR, F1 for each class, ROC--AUC, and PR--AUC. We report mean and sample standard deviation where three training seeds are available. LOGO and fully fine-tuned transformer baselines have one seed and are identified as such; no uncertainty claim is made for them.

\section{Models and Training}

\subsection{Rs-QLoRA detector}

The base checkpoint is \texttt{unsloth/qwen2.5-0.5b-instruct-unsloth-bnb-4bit}, loaded as a two-label \texttt{AutoModelForSequenceClassification} with a maximum sequence length of 512. The frozen base is stored in 4-bit form. Rank-64 adapters target the attention projections \texttt{q}, \texttt{k}, \texttt{v}, and \texttt{o} and the feed-forward \texttt{gate}, \texttt{up}, and \texttt{down} projections. We use $\alpha=128$, zero adapter dropout, no adapter bias, gradient checkpointing, and a trainable sequence-classification head. This configuration has 35,194,624 trainable parameters out of 529,229,184 (6.65\%).

Training uses an effective batch size of 64 (32 per device with two gradient-accumulation steps), up to six epochs, an 8-bit AdamW optimizer, learning rate $5\times10^{-6}$, weight decay 0.03, 6\% warmup, gradient clipping at 1.0, label smoothing 0.01, and cosine scheduling with two restarts. Validation occurs every 100 optimizer steps. Early stopping has patience eight and the checkpoint with greatest validation accuracy is restored. BF16 is used on the H100 GPU. Seeds set the model and training randomness; deterministic cuDNN behavior is requested.

The recorded environment is PyTorch 2.10.0+cu128, Transformers 4.56.2, PEFT 0.19.1, bitsandbytes 0.49.2, Unsloth 2026.6.3, and TRL 0.22.2. Peak memory is measured after resetting PyTorch CUDA peak statistics immediately before training. We distinguish PyTorch allocated memory from process-level memory visible through \texttt{nvidia-smi}.

\subsection{Baselines}

The classical baseline concatenates word 1--2-gram and character-within-word 3--5-gram TF--IDF features and fits logistic regression ($C=4$, 2,000 maximum iterations). The fully fine-tuned baselines are \path{roberta-base} \citep{Liu2019}, \path{microsoft/deberta-v3-base} \citep{He2021}, and \path{answerdotai/ModernBERT-large} \citep{Warner2024}. They use the identical DeepSeek partition, maximum length 512, learning rate $2\times10^{-5}$, effective batch size 32, up to three epochs, 6\% warmup, weight decay 0.01, validation every 200 steps, and early-stopping patience three. Each transformer baseline was run once with seed 1998.

\section{Results}

\subsection{Matched performance is near ceiling}

All four generator-specific detectors achieve high matched accuracy (Table~\ref{tab:matched}). The three-seed standard deviations are small. These results establish in-distribution separability but do not identify what feature family is being used.

\begin{table}[h!]
\centering
\caption{Matched test accuracy across three training seeds (mean $\pm$ sample SD).}
\label{tab:matched}
\begin{tabular}{lc}
\toprule
Training and test generator & Accuracy \\
\midrule
GPT-2 & $98.497 \pm 0.154$\% \\
LFM & $99.776 \pm 0.097$\% \\
DeepSeek & $99.694 \pm 0.057$\% \\
GLM & $99.252 \pm 0.071$\% \\
\bottomrule
\end{tabular}
\end{table}

\subsection{Cross-generator transfer is structured and directional}

Table~\ref{tab:matrix} reports CG recall when a detector trained on one generator is applied to another generator's held-out CG examples. Mean diagonal recall is 99.3\%, whereas mean off-diagonal recall over all four generators is 39.3\%. Restricting off-diagonal cells to LFM, DeepSeek, and GLM raises the mean to 78.4\%.

\begin{table}[h!]
\centering
\caption{Rs-QLoRA CG detection recall (\%) across generators. Rows are training generators and columns are test generators; every cell is averaged over three training seeds. The final column is the detector's mean FPR on its own human test set.}
\label{tab:matrix}
\begin{tabular}{lrrrrr}
\toprule
Train $\backslash$ test & GPT-2 & LFM & DeepSeek & GLM & Human FPR \\
\midrule
GPT-2 & 98.1 & 0.3 & 0.5 & 0.1 & 1.18 \\
LFM & 0.1 & 99.8 & 53.4 & 44.5 & 0.26 \\
DeepSeek & 0.3 & 85.5 & 99.8 & 97.2 & 0.36 \\
GLM & 0.2 & 91.7 & 98.2 & 99.3 & 0.87 \\
\bottomrule
\end{tabular}
\end{table}

Two patterns are robustly visible. First, GPT-2 and the three modern generator domains are nearly disconnected: every cross-boundary recall is 0.5\% or less. A detector that appears nearly perfect on the historical corpus therefore does not detect modern generated reviews, and modern-data detectors do not recover the historical generator.

Second, modern transfer is asymmetric. DeepSeek-trained and GLM-trained detectors recover 85.5--98.2\% of the other modern generators, while the LFM-trained detector recovers 53.4\% of DeepSeek and 44.5\% of GLM. This asymmetry is not explained by source-ID overlap alone: LFM$\rightarrow$GLM and GLM$\rightarrow$LFM have similar source overlap (80.8\% and 81.5\%) but recalls of 44.5\% and 91.7\%, respectively. The result is consistent with a narrower LFM decision boundary and broader signals learned from DeepSeek and GLM, although the current design cannot determine whether those signals originate in model families, provider decoding, or shared generation procedures.

\subsection{Generator-held-out pooled training}

Pooling the other generators substantially improves recall on a held-out modern generator (Table~\ref{tab:logo}). Rs-QLoRA reaches 91.1--96.7\% recall at approximately 2\% human FPR and exceeds TF--IDF in both recall and FPR for all three modern targets. Neither detector transfers to GPT-2. These neural estimates use one seed and their apparent precision should not be read as seed-level certainty.

\begin{table}[h!]
\centering
\caption{LOGO evaluation. The detector is trained on a balanced, budget-matched pool of the other three generators. Held-out OR test texts are excluded from pooled training.}
\label{tab:logo}
\begin{tabular}{lrrrr}
\toprule
Held-out generator & Rs recall & Rs FPR & TF--IDF recall & TF--IDF FPR \\
\midrule
GPT-2 & 0.1\% & 0.48\% & 0.8\% & 1.65\% \\
LFM & 91.1\% & 1.92\% & 84.2\% & 6.40\% \\
DeepSeek & 95.1\% & 2.13\% & 92.3\% & 5.60\% \\
GLM & 96.7\% & 2.01\% & 90.4\% & 6.02\% \\
\bottomrule
\end{tabular}
\end{table}

Class balance also matters operationally. If the observed modern LOGO recall and FPR remained fixed while CG prevalence fell to 1\%, positive predictive value would be only approximately 31--33\% at the uncalibrated 0.5 threshold; at 10\% prevalence it would be approximately 83--84\%. This is a prevalence reweighting, not an additional deployment experiment. It shows why threshold calibration and human review are required in low-prevalence settings.

\subsection{Negative control and matched baselines}

The human-only permutation control obtains 49.31\% accuracy, 48.95\% ROC--AUC, and 49.14\% PR--AUC. Chance performance argues against a general label-channel artifact capable of producing the headline results without a real class signal. It does not eliminate the source-prompt dependence described in Section~3.2.

On the matched DeepSeek split, all neural baselines are near ceiling (Table~\ref{tab:baselines}). DeBERTa-v3-base is best, exceeding mean Rs-QLoRA accuracy by 0.189 percentage points. Rs-QLoRA exceeds RoBERTa-base by 0.183 points and trails ModernBERT-large by 0.082 points. The adapter method's supported advantage is therefore storage and adaptation efficiency, not best matched-distribution accuracy.

\begin{table}[h!]
\centering
\caption{Matched DeepSeek test results. Transformer baselines use one seed; Rs-QLoRA is the mean of three seeds.}
\label{tab:baselines}
\begin{tabular}{lrr}
\toprule
Model & Accuracy & ROC--AUC \\
\midrule
RoBERTa-base & 99.435\% & 99.982\% \\
DeBERTa-v3-base & 99.806\% & 99.996\% \\
ModernBERT-large & 99.700\% & 99.996\% \\
Qwen2.5-0.5B Rs-QLoRA & 99.617\% & 99.960\% \\
\bottomrule
\end{tabular}
\end{table}

\subsection{Error structure on the canonical modern run}

The published DeepSeek seed-1998 adapter makes 24 errors among 5,660 test reviews (99.576\% accuracy): 14 OR reviews are flagged as CG and 10 CG reviews are missed. Errors concentrate in short texts. The error rate is 1.55\% for 8--14 words, 0.64\% for 15--24, 0.33\% for 25--49, 0.16\% for 50--99, and 0.08\% at 100 words or more; the six reviews below eight words contain no error and are too few to interpret. Tools/Home Improvement and Clothing/Shoes/Jewelry have the largest category error rates, both 0.72\%, while all category rates remain below 1\%.

Errors are less confident on average than correct predictions (0.874 versus 0.995 maximum class probability), but 15 of 24 errors still have confidence at least 0.90. False positives often consist of concise, polished statements about value, durability, or use context; false negatives are similarly terse generated statements with product-specific details. Thus, abstention on low confidence would catch some errors but not eliminate high-confidence mistakes. This analysis is descriptive for one run and should not be interpreted as a stable subgroup ranking without additional seeds.

\subsection{Rank stabilization primarily improves convergence efficiency}

The matched three-seed H100 comparison in Table~\ref{tab:efficiency} isolates scaling while holding rank, $\alpha$, data, optimizer, and hardware fixed. Rs-QLoRA improves mean accuracy by only 0.094 percentage points, which is insufficient evidence of an accuracy advantage with three paired seeds. Its clear difference is earlier convergence: mean runtime falls by 45.9\% (1.85$\times$ speedup), and the best validation checkpoint occurs much earlier.

\begin{table}[h!]
\centering
\caption{Rank-stabilized versus standard QLoRA at rank 64 on the matched DeepSeek task (three seeds, one H100 GPU). Memory is the peak PyTorch allocation.}
\label{tab:efficiency}
\begin{tabular}{lrrrr}
\toprule
Scaling & Accuracy, mean $\pm$ SD & Runtime & Best step & Peak allocated \\
\midrule
Rs-QLoRA & $99.617 \pm 0.071$\% & 17.3 min & 800 & 2.05 GiB \\
Standard QLoRA & $99.523 \pm 0.088$\% & 32.0 min & 2,533 & 2.05 GiB \\
\bottomrule
\end{tabular}
\end{table}

Process-level memory was approximately 3.1 GiB in \texttt{nvidia-smi}, including CUDA and process overhead. These values describe adapter training with the stated batch and sequence settings; they are not a general inference-memory benchmark.

\subsection{Rank and scaling ablation on the historical domain}

To test whether the rank-64 choice depends on scaling, we also ran a complete $2\times4$ scaling-by-rank ablation on the historical GPT-2 corpus (Table~\ref{tab:rank}). Every cell uses the same frozen split and three training seeds. Rs scaling exceeds standard scaling at every tested rank, and Rs rank 16 is comparable to standard rank 64. This ablation explains the selected adapter configuration, but it is secondary evidence from the historical domain; it does not override the matched modern comparison in Table~\ref{tab:efficiency}, where both rank-64 variants are at the task ceiling.

\begin{table}[h!]
\centering
\caption{Historical GPT-2 matched accuracy (\%, mean $\pm$ sample SD) by adapter rank and scaling, with three training seeds per cell.}
\label{tab:rank}
\begin{tabular}{rrr}
\toprule
Rank & Standard QLoRA & Rs-QLoRA \\
\midrule
8 & $96.008 \pm 0.221$ & $97.315 \pm 0.338$ \\
16 & $97.044 \pm 0.080$ & $97.868 \pm 0.082$ \\
32 & $97.297 \pm 0.383$ & $98.139 \pm 0.168$ \\
64 & $97.792 \pm 0.093$ & $98.740 \pm 0.089$ \\
\bottomrule
\end{tabular}
\end{table}

\section{Discussion}

\subsection{What the shared modern signal does and does not mean}

The 78.4\% modern-only off-diagonal recall and 91.1--96.7\% modern LOGO recall are substantially above the near-zero GPT-2 transfer. This is useful evidence that the tested modern corpora contain shared, exploitable regularities beyond one detector's matched generator. The directional structure adds information: broad DeepSeek/GLM boundaries and a narrow LFM boundary would not arise from a single scalar notion of ``AI-likeness.'' A practical ensemble or pooled detector can use this diversity more effectively than a detector trained on one generator.

The finding nevertheless admits several causal explanations. Modern corpora share source reviews, prompt wording, length control, acceptance filtering, normalization, and broadly similar instruction-following generation. Their common signal may reflect contemporary model behavior, the common generation pipeline, source-conditioned paraphrasing, or a mixture of these. The 69--82\% pairwise and approximately 94\% pooled source overlap means that high LOGO recall is not evidence for novel products or prompts. A future factorial experiment should cross generators with independently held-out source groups and matched decoding settings. Only then can model-family, prompt, and source effects be estimated separately.

The GPT-2 discontinuity is equally important. It demonstrates that ``machine generated'' is not one stationary distribution. A detector trained on one technological era can fail almost completely on another. The present results support continuous data refresh and multi-generator training rather than permanent deployment of a single historical detector.

\subsection{Operational interpretation}

At roughly 2\% FPR, a modern LOGO detector may be useful for prioritizing review moderation within a similar English product-review pipeline. It should not be used as an autonomous accusation mechanism. Even a stable FPR produces many false alerts when generated reviews are rare, and the study does not test paraphrasing attacks, human editing, multilingual text, different review platforms, or arbitrary future models. Scores require deployment-specific calibration, monitoring for generator drift, and human adjudication.

The transfer matrix also does not perform generator attribution. Low or high response patterns suggest generator fingerprints, but a binary classifier trained against OR data is not an identified multiclass or open-set attribution model. Attribution would require matched prompts and decoding across generators, source-grouped evaluation, unknown-generator rejection, and direct multiclass training.

\section{Limitations and Ethical Considerations}

The study has six primary limitations. First, all content is English and restricted to product reviews from one source population. Second, modern generator names are provider request aliases; immutable provider snapshots and the exact local LFM serving build were not captured. The committed raw outputs are reproducible artifacts, but byte-identical regeneration is impossible because decoding was unseeded. Third, generator, provider, prompting, and sampling effects are confounded. Fourth, the row-level and cross-generator designs do not jointly hold out grounding-source identity. Fifth, LOGO and fully fine-tuned transformer results use one training seed. Sixth, no adversarial paraphrase, human-editing, multilingual, temporal-drift, or cross-platform evaluation was performed.

Generated-review detection can support moderation, measurement, and provenance research, but false positives can harm legitimate reviewers. The model output should be treated as a risk score rather than proof of deception or authorship. The released corpora contain generated reviews grounded on public review text; downstream users remain responsible for the upstream data terms and for avoiding uses that facilitate review manipulation.

\section{Conclusion}

Near-ceiling matched accuracy concealed a much more informative transfer structure. GPT-2 and the tested modern generators form almost disjoint detection domains, while LFM, DeepSeek, and GLM share transferable but directional signals under a common grounded generation pipeline. Pooled training detects a held-out modern generator with 91.1--96.7\% recall at about 2\% human FPR, but extensive source-prompt overlap prevents interpreting this as fully out-of-distribution generalization. Rs-QLoRA does not beat the strongest fully fine-tuned matched baseline; its demonstrated benefit is efficient adaptation, including a 45.9\% training-time reduction relative to standard QLoRA at comparable accuracy. The resulting claim is deliberately bounded: compact detectors can exploit shared signals across the studied modern review generators, but robust machine-authorship detection requires generator-diverse, source-grouped, and temporally refreshed evaluation.

\section*{Conflict of Interest Statement}
The authors declare that the research was conducted in the absence of any commercial or financial relationships that could be construed as a potential conflict of interest.

\section*{Author Contributions}
ZR conceived the study, designed and implemented the experiments, analyzed the results, and prepared the manuscript. SEF supervised the research and reviewed the manuscript. RF co-supervised the work and provided critical feedback. All authors reviewed and approved the final manuscript.

\section*{Funding}
This research received no specific grant from any funding agency in the public, commercial, or not-for-profit sectors.

\section*{Acknowledgments}
The authors acknowledge the creators and maintainers of the public source dataset and open-source software used in this study.

\section*{Generative AI Statement}
Generative AI tools were used during manuscript revision for drafting support, language editing, and code-assisted reconciliation of experimental records. The authors reviewed the underlying artifacts, verified the reported analyses and citations, determined the scientific interpretation, and take responsibility for the final text. This statement should be reconciled with the journal's policy and disclosure form before submission.

\section*{Data Availability Statement}
The modern datasets are available on Hugging Face for \href{https://huggingface.co/datasets/Flowerly/modern-fake-reviews-lfm}{LFM}, \href{https://huggingface.co/datasets/Flowerly/modern-fake-reviews}{DeepSeek}, and \href{https://huggingface.co/datasets/Flowerly/modern-fake-reviews-glm}{GLM}. The \href{https://www.kaggle.com/datasets/mexwell/fake-reviews-dataset}{historical source dataset} is available on Kaggle. Code, structured metrics, run manifests, and dataset provenance are available in the \href{https://github.com/PurplexyFlower/Fake-review}{project repository}. Published adapter weights and their evidence manifest are available in the \href{https://huggingface.co/KandirResearch/fake-review-rs-qlora-adapters}{adapter repository}.

\bibliographystyle{Frontiers-Harvard}
\bibliography{references}

\end{document}
